\documentclass[11pt,a4paper,UTF8]{ctexart}
\usepackage{geometry}
\geometry{left=18mm,right=18mm,top=24mm,bottom=24mm,headheight=22pt,footskip=12mm}
\usepackage{amsmath,amssymb,mathtools,bm,esint,extarrows,centernot}
\usepackage{xcolor}
\usepackage{tcolorbox}
\tcbuselibrary{skins,breakable}
\usepackage{fancyhdr}
\usepackage{titlesec}
\usepackage{hyperref}
\usepackage{tikz}
\usepackage{booktabs}
\usepackage{tabularx}

% --- Color Palette (Purple & DeepPink Academic Theme) ---
\definecolor{DeepPurple}{HTML}{4C1D95}
\definecolor{TitlePurple}{HTML}{581C87}
\definecolor{SubPurple}{HTML}{6B21A8}
\definecolor{DeepPink}{HTML}{FF1493}
\definecolor{LightPinkBg}{HTML}{FFF5F8}
\definecolor{LightPurpleBg}{HTML}{F8F5FF}
\definecolor{GoldAccent}{HTML}{D97706}
\definecolor{DarkText}{HTML}{1F2937}

\hypersetup{
    colorlinks=true,
    linkcolor=TitlePurple,
    citecolor=DeepPink,
    urlcolor=SubPurple,
    pdfborder={0 0 0}
}

% --- Typography & Spacing ---
\linespread{1.3}
\setlength{\parskip}{0.35em plus 0.1em minus 0.1em}
\setlength{\parindent}{0pt}

% --- Section Titles Styling ---
\titleformat{\section}
  {\Large\bfseries\color{TitlePurple}}{\thesection}{1em}{}
  [\vspace{1mm}{\color{TitlePurple!30}\hrule height 1pt}]
\titleformat{\subsection}
  {\large\bfseries\color{SubPurple}}{\thesubsection}{0.8em}{}
\titleformat{\subsubsection}
  {\normalsize\bfseries\color{DeepPurple}}{\thesubsubsection}{0.6em}{}

% --- tcolorbox Environments ---
% 1. Theorem / Formula / Method / Analysis Box (Deep Purple)
\newtcolorbox{thmbox}[1][]{
    enhanced,
    breakable,
    colback=LightPurpleBg,
    colframe=TitlePurple,
    arc=3pt,
    boxrule=0pt,
    leftrule=3.5pt,
    left=10pt,
    right=10pt,
    top=8pt,
    bottom=8pt,
    title=#1,
    coltitle=white,
    colbacktitle=TitlePurple,
    fonttitle=\bfseries\small,
    attach boxed title to top left={yshift=-2mm, xshift=4mm},
    boxed title style={boxrule=0pt, arc=2pt}
}

% 2. Solution / Formula / Derivation Box (DeepPink)
\newtcolorbox{solbox}[1][]{
    enhanced,
    breakable,
    colback=LightPinkBg,
    colframe=DeepPink,
    arc=3pt,
    boxrule=0pt,
    leftrule=3.5pt,
    left=10pt,
    right=10pt,
    top=8pt,
    bottom=8pt,
    title=#1,
    coltitle=white,
    colbacktitle=DeepPink,
    fonttitle=\bfseries\small,
    attach boxed title to top left={yshift=-2mm, xshift=4mm},
    boxed title style={boxrule=0pt, arc=2pt}
}

% 3. Learning Goals / Summary Box
\newtcolorbox{targetbox}[1][]{
    enhanced,
    breakable,
    colback=LightPurpleBg,
    colframe=DeepPurple,
    arc=3pt,
    boxrule=1pt,
    left=10pt,
    right=10pt,
    top=8pt,
    bottom=8pt,
    title=#1,
    coltitle=white,
    colbacktitle=DeepPurple,
    fonttitle=\bfseries\small,
    attach boxed title to top left={yshift=-2mm, xshift=4mm},
    boxed title style={boxrule=0pt, arc=2pt}
}

\begin{document}

% ------------------- 讲义封面 (Cover Page) -------------------
\begin{titlepage}
\thispagestyle{empty}
\begin{tikzpicture}[remember picture,overlay]
    \fill[DeepPurple] (current page.north west) rectangle ([yshift=-7cm]current page.north east);
    \draw[DeepPink, line width=3.5pt] ([yshift=-7cm]current page.north west) -- ([yshift=-7cm]current page.north east);
    \draw[GoldAccent, line width=1pt] ([yshift=-7.12cm]current page.north west) -- ([yshift=-7.12cm]current page.north east);
    
    \node[anchor=north east, opacity=0.12, text=white] at ([xshift=-1.5cm, yshift=-0.5cm]current page.north east) {
        \fontsize{70}{70}\selectfont\bfseries 2027
    };
\end{tikzpicture}

\vspace*{0.8cm}
\begin{center}
    {\color{white}\fontsize{26}{32}\selectfont\bfseries 2027 考研数学(二) 核心讲义}\\[0.6cm]
    {\color{white}\fontsize{13}{18}\selectfont\bfseries 全国硕士研究生招生考试 · 统考数学(二) 核心精讲全书}\\[1.5cm]
\end{center}

\vspace{3.5cm}
\begin{center}
    {\color{GoldAccent}\large\bfseries $\blacktriangleright$ 基础强化 · 核心题解 · 考点深水区 $\blacktriangleleft$}\\[0.8cm]
    {\color{TitlePurple}\fontsize{24}{28}\selectfont\bfseries 第七章 一元函数微分学的应用(物理应用和经济应用)}\\[0.6cm]
    {\color{DeepPink}\large\bfseries —— 体系化理论精讲与公式全景 ——}\\[1.5cm]
\end{center}

\vfill

\begin{center}
\begin{tcolorbox}[
    colback=white,
    colframe=DeepPurple,
    width=0.88\textwidth,
    arc=4pt,
    boxrule=1.2pt,
    halign=center,
    drop shadow=gray!30
]
    \vspace{3mm}
    {\bfseries\large 编著团队：EscoffierZhou}\\[2.5mm]
    {\bfseries\color{TitlePurple} 目标院校：中国科学院大学 (UCAS) 计算机专硕 (22408)}\\[2.5mm]
    {\small\color{gray} 备考铁律：手算真题数字 · 错题白纸重做 · 408 客观题为王}\\[2.5mm]
    {\small 编制版本：2027 届首发版 · \today}
    \vspace{2mm}
\end{tcolorbox}
\end{center}
\vspace{0.8cm}
\end{titlepage}

% ------------------- 目录与页面主体配置 -------------------
\pagestyle{fancy}
\fancyhf{}
\fancyhead[L]{\small\color{gray}2027 考研数学(二) · 核心讲义系列}
\fancyhead[R]{\small\color{TitlePurple}\nouppercase{\leftmark}}
\fancyfoot[C]{\small\color{gray}— 第 \thepage\ 页 —}
\fancyfoot[R]{\small\color{gray}UCAS 22408}
\renewcommand{\headrulewidth}{0.5pt}

\pagenumbering{Roman}
\tableofcontents
\newpage
\pagenumbering{arabic}


\begin{targetbox}[本章复习目标与核心知识图谱]
\begin{itemize}
  \item 目的[1]:掌握导数的物理本质——变化率模型与相关变化率(数一/数二高频考点)
\end{itemize}
\end{targetbox}

\begin{targetbox}[本章复习目标与核心知识图谱]
\begin{itemize}
  \item 目的[2]:掌握质点直线运动中的位移、速度、加速度与高阶导数链式微分关系
\end{itemize}
\end{targetbox}

\begin{targetbox}[本章复习目标与核心知识图谱]
\begin{itemize}
  \item 目的[3]:掌握平面曲线运动中的速度分解、法向/切向加速度与曲率动力学模型
\end{itemize}
\end{targetbox}

\begin{targetbox}[本章复习目标与核心知识图谱]
\begin{itemize}
  \item 目的[4]:熟练建立物理量之间的几何/代数约束方程，并结合隐函数求导法与链式法则求解变化率
\end{itemize}
\end{targetbox}

\begin{targetbox}[本章复习目标与核心知识图谱]
\begin{itemize}
  \item 目的[5]:掌握导数在经济学中的边际分析(边际成本、边际收益、边际利润及利润最大化准则)
\end{itemize}
\end{targetbox}

\begin{targetbox}[本章复习目标与核心知识图谱]
\begin{itemize}
  \item 目的[6]:掌握需求价格弹性公式、弹性与收益的关系，熟练运用导数分析经济决策
\end{itemize}
\end{targetbox}


\section{1.导数的物理本质:变化率与相关变化率}


\begin{equation*}
\frac{\mathrm{d}y}{\mathrm{d}t} = \lim_{\Delta t \to 0} \frac{\Delta y}{\Delta t}
\end{equation*}
\begin{thmbox}[原理[1]: 变化率的物理定义]

若物理量$y$依赖于时间参数$t$，则$y$对$t$的导数即为$y$随时间$t$变化的\textbf{瞬时变化率}：


- 若$\frac{\mathrm{d}y}{\mathrm{d}t} > 0$，表示物理量随时间增加(如膨胀、上升、升温、加速)；
- 若$\frac{\mathrm{d}y}{\mathrm{d}t} < 0$，表示物理量随时间减少(如收缩、下降、降温、减速)。

\end{thmbox}

\begin{equation*}
F(x, y, \dots) = 0
\end{equation*}
\begin{thmbox}[原理[2]: 相关变化率模型与标准SOP流程]

设两个或多个物理量$x, y, \theta$都是时间$t$的可导函数，它们之间通过某种几何约束关系或物理规律满足关联方程：


若已知其中某些量的变化率(如$\frac{\mathrm{d}x}{\mathrm{d}t}$)，求另一物理量对时间的变化率$\frac{\mathrm{d}y}{\mathrm{d}t}$。

[相关变化率标准解题SOP(步控流程)]

操作[1]:选元定标：设出所有随时间变化的变量(记为$x(t), y(t), \theta(t)$等)。

注意:变化过程中的动量严禁在求导前代入固定数值！

操作[2]:几何/物理建模：利用勾股定理、相似三角形、三角函数定义或体积公式建立几何约束方程$F(x, y) = 0$。

操作[3]:对时间$t$全求导：方程两边关于时间$t$求导(视各变量为$t$的复合函数，严格应用链式法则)。

操作[4]:代入特定瞬时值：仅在求导完成后，将题设中“特定瞬间”的几何位置参数和已知变化率代入，解出目标变化率。

\end{thmbox}

\section{2.质点直线运动与微分关系}

\begin{thmbox}[原理[1]: 质点直线运动中的高阶导数链]

设质点沿直线运动的位置坐标随时间的变化规律为位移函数$s = s(t)$：

- \textbf{瞬时速度(一阶导数):} $v(t) = \frac{\mathrm{d}s}{\mathrm{d}t} = s'(t)$
- \textbf{瞬时加速度(二阶导数):} $a(t) = \frac{\mathrm{d}v}{\mathrm{d}t} = \frac{\mathrm{d}^2s}{\mathrm{d}t^2} = s''(t)$
- \textbf{加加速度/变加速度(三阶导数Jerk):} $j(t) = \frac{\mathrm{d}a}{\mathrm{d}t} = \frac{\mathrm{d}^3s}{\mathrm{d}t^3} = s'''(t)$

\end{thmbox}

\begin{equation*}
a = \frac{\mathrm{d}v}{\mathrm{d}t} = \frac{\mathrm{d}v}{\mathrm{d}s} \cdot \frac{\mathrm{d}s}{\mathrm{d}t} = v \frac{\mathrm{d}v}{\mathrm{d}s} = \frac{\mathrm{d}}{\mathrm{d}s}\left(\frac{1}{2}v^2\right)
\end{equation*}
\begin{thmbox}[原理[2]: 以位移为自变量的链式法则变换]

在动力学问题中，加速度或外力常由位置决定(如弹簧弹力、万有引力$a = a(s)$)。此时时间参数$t$被消去，利用复合求导链式法则进行转化：


本质作用:此关系是建立牛顿第二定律微分方程的桥梁：$F = ma \implies F(s) = mv \frac{\mathrm{d}v}{\mathrm{d}s}$，分离变量即可直接积分求解速度与位移的关系。

\end{thmbox}

\section{3.平面曲线运动与切向/法向加速度(结合曲率)}


\begin{equation*}
\vec{a} = a_\tau \vec{\tau} + a_n \vec{n}
\end{equation*}

\begin{equation*}
a_\tau = \frac{\mathrm{d}v}{\mathrm{d}t} = \frac{\mathrm{d}^2s}{\mathrm{d}t^2}
\end{equation*}

\begin{equation*}
a_n = \frac{v^2}{R} = K v^2
\end{equation*}

\begin{equation*}
a = \vert{}\vec{a}\vert{} = \sqrt{a_\tau^2 + a_n^2} = \sqrt{\left(\frac{\mathrm{d}v}{\mathrm{d}t}\right)^2 + (K v^2)^2}
\end{equation*}
\begin{thmbox}[原理[1]: 切向加速度与法向加速度的正交分解]

当质点沿平面曲线$y = f(x)$运动，弧长参数为$s$，瞬时速率为$v = \frac{\mathrm{d}s}{\mathrm{d}t}$时，其加速度向量$\vec{a}$可正交分解为切向分量$\vec{a}_\tau$与法向分量$\vec{a}_n$：


- \textit{}切向加速度$a_\tau$:\textbf{ 反映}速率大小改变的快慢\textit{}：


- \textit{}法向加速度$a_n$:\textbf{ 反映}运动速度方向改变的快慢\textit{}，与轨迹在当前点的曲率$K$及曲率半径$R$直接绑定：


- \textbf{全加速度大小:} 


\end{thmbox}

\begin{equation*}
K = \frac{\vert{}y''\vert{}}{[1 + (y')^2]^{\frac{3}{2}}}
\end{equation*}
\begin{thmbox}[原理[2]: 平面直角坐标下的曲率计算]

若曲线方程为$y = f(x)$，在点$(x, y)$处的曲率$K$为：


曲率半径为$R = \frac{1}{K}$。

\end{thmbox}

\section{5.经典强化例题精解}


\subsection{例题 1: 相关变化率(仰角与高度模型)}

观测站位于离火箭发射台水平距离为$3\text{ km}$的平地上。若火箭垂直向上发射，当火箭高度达到$4\text{ km}$时，其瞬时上升速度为$800\text{ km/h}$。求此时观测站仪器对火箭观测仰角的变化率。

\textbf{\color{DeepPurple}【思路点拨】} 利用直角三角形建立正切函数几何约束方程$y = 3\tan\theta$，两端对时间$t$求导应用链式法则，代入特定瞬间的几何量解出仰角变化率


\begin{solbox}[分步推导与答题规范]
\textbf{[1]} 选元与几何建模：

设在时刻$t$，火箭垂直高度为$y(t)$，观测站仪器的仰角为$\theta(t)$。

由直角三角形三角函数关系：


\begin{equation*}
\tan\theta = \frac{y}{3} \implies y = 3\tan\theta
\end{equation*}


\textbf{[2]} 两端关于时间$t$求导：

方程两端关于时间$t$求导(视$y$与$\theta$为$t$的复合函数)：


\begin{equation*}
\frac{\mathrm{d}y}{\mathrm{d}t} = 3\sec^2\theta \cdot \frac{\mathrm{d}\theta}{\mathrm{d}t}
\end{equation*}


\textbf{[3]} 确定特定瞬间的几何量：

当火箭高度$y = 4\text{ km}$时：

- 水平距离为$3\text{ km}$，斜边视线距离为$\sqrt{3^2 + 4^2} = 5\text{ km}$；
- 此时$\cos\theta = \frac{3}{5} \implies \sec\theta = \frac{5}{3}$；
- 瞬时上升速度已知：$\frac{\mathrm{d}y}{\mathrm{d}t} = 800\text{ km/h}$。

\textbf{[4]} 代入计算仰角变化率：


\begin{equation*}
800 = 3 \cdot \left(\frac{5}{3}\right)^2 \cdot \frac{\mathrm{d}\theta}{\mathrm{d}t} = \frac{25}{3} \cdot \frac{\mathrm{d}\theta}{\mathrm{d}t}
\end{equation*}


解得：


\begin{equation*}
\frac{\mathrm{d}\theta}{\mathrm{d}t} = \frac{800 \times 3}{25} = 96\text{ rad/h}
\end{equation*}

\end{solbox}


\subsection{例题 2: 相关变化率(倒圆锥漏水模型)}

一倒置圆锥形水漏斗，其底面圆半径为$R = 6\text{ cm}$，高为$H = 12\text{ cm}$。水以$2\text{ cm}^3\text{/s}$的恒定速率从底部小孔漏出。当水深为$4\text{ cm}$时，求漏斗中水面下降的速率。

\textbf{\color{DeepPurple}【思路点拨】} 利用相似三角形将半径$r$用高$h$表示以消元，建立单一变量体积公式$V(h)$，两边对$t$求导后代入瞬时水深求解


\begin{solbox}[分步推导与答题规范]
\textbf{[1]} 选元与几何相似消元：

设在时刻$t$，水深为$h(t)$，对应水面半径为$r(t)$，水的体积为$V(t)$。

由圆锥截面的相似三角形性质：


\begin{equation*}
\frac{r}{h} = \frac{R}{H} = \frac{6}{12} = \frac{1}{2} \implies r = \frac{1}{2}h
\end{equation*}


\textbf{[2]} 建立单变量体积函数：


\begin{equation*}
V = \frac{1}{3}\pi r^2 h = \frac{1}{3}\pi \left(\frac{1}{2}h\right)^2 h = \frac{1}{12}\pi h^3
\end{equation*}


\textbf{[3]} 两端对时间$t$求导：


\begin{equation*}
\frac{\mathrm{d}V}{\mathrm{d}t} = \frac{1}{12}\pi \cdot 3h^2 \frac{\mathrm{d}h}{\mathrm{d}t} = \frac{1}{4}\pi h^2 \frac{\mathrm{d}h}{\mathrm{d}t}
\end{equation*}


\textbf{[4]} 代入已知瞬时值计算：

漏水说明体积随时间减少，$\frac{\mathrm{d}V}{\mathrm{d}t} = -2\text{ cm}^3\text{/s}$，水深$h = 4\text{ cm}$：


\begin{equation*}
-2 = \frac{1}{4}\pi (4)^2 \frac{\mathrm{d}h}{\mathrm{d}t} = 4\pi \frac{\mathrm{d}h}{\mathrm{d}t} \implies \frac{\mathrm{d}h}{\mathrm{d}t} = -\frac{1}{2\pi}\text{ cm/s}
\end{equation*}


故水面下降的速率为$\frac{1}{2\pi}\text{ cm/s}$。
\end{solbox}


\subsection{例题 3: 相关变化率(滑梯靠墙下滑模型)}

长为$5\text{ m}$的梯子靠在垂直的墙上。底端以$2\text{ m/s}$的匀速沿水平地面向远离墙面方向滑动。当底端离墙面$3\text{ m}$时，求顶端沿墙面下滑的速度，以及梯子与地面夹角$\theta$的变化率。

\textbf{\color{DeepPurple}【思路点拨】} 建立勾股定理约束方程$x^2 + y^2 = 25$对时间求导求解顶端下滑速度；通过余弦定义$\cos\theta = \frac{x}{5}$对时间求导求解角速度


\begin{solbox}[分步推导与答题规范]
\textbf{[1]} 勾股定理建模与求导：

设底端离墙距离为$x(t)$，顶端离地高度为$y(t)$。满足几何约束：


\begin{equation*}
x^2 + y^2 = 5^2 = 25
\end{equation*}


两端对时间$t$求导：


\begin{equation*}
2x\frac{\mathrm{d}x}{\mathrm{d}t} + 2y\frac{\mathrm{d}y}{\mathrm{d}t} = 0 \implies x\frac{\mathrm{d}x}{\mathrm{d}t} + y\frac{\mathrm{d}y}{\mathrm{d}t} = 0
\end{equation*}


\textbf{[2]} 计算特定瞬时的顶端速度：

当$x = 3\text{ m}$时，代入勾股方程得$y = \sqrt{25 - 3^2} = 4\text{ m}$。

已知$\frac{\mathrm{d}x}{\mathrm{d}t} = 2\text{ m/s}$，代入得：


\begin{equation*}
3(2) + 4\frac{\mathrm{d}y}{\mathrm{d}t} = 0 \implies \frac{\mathrm{d}y}{\mathrm{d}t} = -\frac{3}{2}\text{ m/s}
\end{equation*}


说明顶端沿墙面下滑的速度大小为$1.5\text{ m/s}$。

\textbf{[3]} 求夹角$\theta$的变化率：

由几何关系$\cos\theta = \frac{x}{5}$，两端对$t$求导：


\begin{equation*}
-\sin\theta \frac{\mathrm{d}\theta}{\mathrm{d}t} = \frac{1}{5}\frac{\mathrm{d}x}{\mathrm{d}t}
\end{equation*}


此时$\sin\theta = \frac{y}{5} = \frac{4}{5}$，代入得：


\begin{equation*}
-\frac{4}{5} \frac{\mathrm{d}\theta}{\mathrm{d}t} = \frac{1}{5} \times 2 \implies \frac{\mathrm{d}\theta}{\mathrm{d}t} = -\frac{1}{2}\text{ rad/s}
\end{equation*}

\end{solbox}


\subsection{例题 4: 曲线运动(法向加速度与曲率结合)}

质点沿抛物线轨道$y = \frac{1}{2}x^2$运动。当质点到达点$(1, \frac{1}{2})$时，已知质点在该点的瞬时速率为$v = 2\sqrt{5}\text{ m/s}$，且速率随时间的增加率为$\frac{\mathrm{d}v}{\mathrm{d}t} = 3\text{ m/s}^2$。求质点在该瞬时的全加速度大小。

\textbf{\color{DeepPurple}【思路点拨】} 切向加速度直接由速率导数给出；法向加速度通过轨迹曲率公式$K$及$a_n = Kv^2$计算；最后合成全加速度


\begin{solbox}[分步推导与答题规范]
\textbf{[1]} 切向加速度计算：

切向加速度反映速率大小改变的快慢，直接由已知条件给出：


\begin{equation*}
a_\tau = \frac{\mathrm{d}v}{\mathrm{d}t} = 3\text{ m/s}^2
\end{equation*}


\textbf{[2]} 计算轨迹在点$x = 1$处的曲率$K$：

- $y' = x \implies y'(1) = 1$；
- $y'' = 1 \implies y''(1) = 1$；
- 代入平面曲线曲率公式：

 
\begin{equation*}
K = \frac{\vert{}y''\vert{}}{[1 + (y')^2]^{\frac{3}{2}}} = \frac{1}{[1 + 1^2]^{\frac{3}{2}}} = \frac{1}{2^{\frac{3}{2}}} = \frac{1}{2\sqrt{2}}
\end{equation*}


\textbf{[3]} 计算法向加速度$a_n$：


\begin{equation*}
a_n = K v^2 = \frac{1}{2\sqrt{2}} \times (2\sqrt{5})^2 = \frac{20}{2\sqrt{2}} = 5\sqrt{2}\text{ m/s}^2
\end{equation*}


\textbf{[4]} 合成全加速度大小：


\begin{equation*}
a = \sqrt{a_\tau^2 + a_n^2} = \sqrt{3^2 + (5\sqrt{2})^2} = \sqrt{9 + 50} = \sqrt{59}\text{ m/s}^2
\end{equation*}

\end{solbox}

\end{document}
